% LỜI TỰA --- Quyển XV
\chapter*{Lời Tựa}
\addcontentsline{toc}{chapter}{Lời Tựa}
\markboth{Lời Tựa}{Lời Tựa}

\begin{truyen}{Lời Tựa}{What Teachers Want to Tell Students}
\chuhoa{G}{iáo viên bắt đầu bộ sách này với một câu hỏi đơn giản: làm thế nào để mỗi buổi học có thể để lại thứ gì đó đáng nhớ hơn là bài thi?}
\textit{(A teacher began this series with a simple question: how can each class leave behind something more memorable than a test?)}

Có những điều thầy muốn nói với học trò~--- nhưng tiết học hết quá nhanh, và không phải lúc nào cũng đúng lúc.
\textit{(There are things teachers want to say to students~--- but class ends too quickly, and the timing is never quite right.)}

Quyển mười lăm chọn một giọng khác: không lên lớp quá cao, không làm quá nhiều khẩu hiệu. Mỗi bài ngắn như một lời thầy nói lại một điều quen thuộc nhưng cố nói gần hơn, thật hơn, và có ích hơn. Nếu em đang đi học, thầy mong em đọc quyển này như đang ngồi nghe một người lớn muốn nói chuyện đàng hoàng với mình.
\textit{(Volume fifteen chooses a different voice: not lecturing from too high above, not filling pages with empty slogans. Each short piece is like a teacher saying something familiar again but closer, truer, and more useful. If you are a student, read this as though sitting with an adult who wants to speak properly with you.)}

Bộ sách <<Truyện Tích Cảm Hứng Song Ngữ>> gồm 24 quyển, mỗi quyển một chủ đề riêng. Quyển XV này mang chủ đề: \textbf{thầy nói thẳng, gần và chậm về học tập, tính cách, đường đời}.
\textit{(The series ``Inspiring Bilingual Tales'' contains 24 volumes, each with its own theme. Volume XV focuses on: \textbf{thầy nói thẳng, gần và chậm về học tập, tính cách, đường đời}.})
\end{truyen}

\begin{baihoc}
Lời thầy gần và thật hơn một nghìn lần so với lời thầy sáo rỗng~--- học sinh biết ngay sự khác biệt.
\textit{(A teacher's close and honest word is worth a thousand empty ones~--- students always know the difference.)}
\end{baihoc}

\begin{ghinhoanh}
\textbf{Short English to remember:}

What a teacher truly says lives long after the lesson ends.

Speak close to the student, not above them.
\end{ghinhoanh}

\begin{tuhocnhanh}
1. Điều thầy nói nào trong quyển này em ước mình được nghe từ sớm hơn? \textit{(What did a teacher say in this volume that you wish you'd heard earlier?)}

2. Em đã từng thực sự lắng nghe lời thầy cô chưa? Lần nào đáng nhớ nhất? \textit{(Have you ever truly listened to a teacher? What was the most memorable time?)}

3. Nếu em có thể nói một điều thật lòng với thầy cô, đó là điều gì? \textit{(If you could say one honest thing to a teacher, what would it be?)}
\end{tuhocnhanh}

\begin{center}
\textit{\color{softcopper}``Thầy nói không phải để em ghi nhớ công thức~--- mà để em sống tốt hơn.''}
\end{center}

\begin{flushright}
\textbf{Tôi Nguyễn Văn Sang}\\
\textit{GV THPT Nguyễn Hữu Cảnh - TP HCM}
\end{flushright}

\ngancach
